\documentclass[../../main.tex]{subfiles}
\begin{document}

\begin{center}
  {\fontsize{48pt}{58pt}\selectfont \textbf{2003}\normalsize\textbf{年度}}
\end{center}

\begin{table}[htb]
  \centering
  \caption{出題テーマと難易度}
  \begin{tabular}{|c|p{9cm}|c|}
    \hline
    \textbf{問題番号} & \textbf{テーマ} & \textbf{難易度} \\
    \hline
    第1問 & 対数差で定まる数列の和 & \\
    \hline
    第2問 & 正弦曲線の法線で囲む回転体 & \\
    \hline
    第3問 & 垂直条件と等面積の四面体 & \\
    \hline
    第4問 & 多項式の割り切りと三乗根 & \\
    \hline
    第5問 & 行列集合の逆行列の存在条件 & \\
    \hline
    第6問 & リーグ戦の勝敗確率 & \\
    \hline
  \end{tabular}
\end{table}

\begin{table}[htb]
  \centering
  \caption{解答}
  \renewcommand{\arraystretch}{1.4}
  \begin{tabular}{|c|c|p{9cm}|}
    \hline
    \textbf{問題番号} & \textbf{小問} & \textbf{答え} \\
    \hline
    第1問 & -- & $\sum_{k=1}^n a_k=\dfrac{2n}{n+1}$ \\
    \hline
    第2問 & -- & $V=\dfrac{13\pi^4}{48}-\dfrac{\pi^2}{8}$ \\
    \hline
    第3問 & -- & 四面体は正四面体（証明問題） \\
    \hline
    第4問 & -- & $x^2+x+1$で割り切れる \\
    \hline
    第5問 & -- &  \\
    \hline
    第6問 & -- & $n(n-1)(n-2)2^{-(2n-3)}\left(1-2^{-(n-3)}\right)$ \\
    \hline
  \end{tabular}
\end{table}

\paragraph{試験形態}
6問．試験時間は要確認．

\end{document}
